\documentclass[../../main.tex]{subfiles}

\begin{document}

\section{Obbiettivi attuali} \label{sec:obbiettivi-attuali}
  %    - propedeuticità al riempimento dei caf con roba vera
  %    - validazione input
  %    - testare analisi dati su oggetti completi
  Al momento della stesura di questa tesi fake-reco adempie ai seguenti
  compiti:

  \paragraph{Propedeuticità}
    I CAF di DUNE sono delle strutture dati nuove e ancora in via di sviluppo
    che i ricercatori devono imparare a usare, riempirli con la verità Monte
    Carlo permette di studiarli a fondo ed eventualmente di individuare
    storture nel loro design.

  \paragraph{Validazione dell'input}
    La maggior parte dei \textit{tool} di analisi dati sarà sviluppata in
    maniera naturale per lavorare sui CAF; avere un layer di traduzione
    uno a uno tra l'output di EDep-Sim e i CAF elimina la necessità di creare
    strumenti appositi per i dati simulati e, nel malaugurato caso in cui
    lo standard dei formati in ingresso dovessero cambiare, riduce lo
    sforzo reimplementativo alla sola fake-reco.

  \paragraph{Test degli strumenti di analisi}
    Per poter testare i workflow di analisi dati è necessario avere i dati
    su cui lavorare.
    Delle strutture dati riempite con informazioni realistiche seppur troppo
    perfette permettono di testare più a fondo gli strumenti di
    analisi rispetto a strutture riempite con dati nulli o completamente
    casuali.

\end{document}